\documentclass[12pt]{article}
\usepackage{amssymb,amsmath}
\usepackage{anyfontsize}
\thispagestyle{empty}
\usepackage{tikz}
\usetikzlibrary{calc}
\usepackage{setspace}
\usepackage{graphicx}
\usepackage[margin=0.65in]{geometry}
\usepackage{float}
\usepackage[letterspace=150]{microtype}


\begin{document}

%Border

\begin{tikzpicture}[overlay,remember picture]
\draw [line width=3pt,rounded corners=0pt,]
    ($ (current page.north west) + (1.25cm,-1.25cm) $)
    rectangle
    ($ (current page.south east) + (-1.25cm,1.25cm) $);       
\end{tikzpicture}
\begin{center}\includegraphics[scale=.35]{RUMAlogo.png}\\
presents... \\
\begin{spacing}{1}
{\fontsize{40}{44}\selectfont  \textsc{
Symplectic Geometry}} \end{spacing}
 

~~\\
\begin{spacing}{1.5}
{\fontsize{24}{24} \selectfont A Lecture by Professor Christopher
  Woodward}  \end{spacing} 
\large Dept. of Mathematics, Rutgers University \\

\vspace{5mm}
\normalsize
\textsc{Abstract:}

\Large 
Symplectic geometry arose as a mathematical framework 
for Hamiltonian dynamics in physics and is now one of the most
active areas of research in mathematics, especially focused on the
``mirror symmetry
conjectures'' relating symplectic and algebraic geometry and the ``Arnold conjectures''
on time-dependent Hamiltonian flows.

\vspace{2mm}
I'll give a quick introduction to symplectic geometry based on the following 
real life problem:  Put a rubber band around your favorite math book 
and throw it in the air, with the cover facing you, and try to catch
it after it's done a 360 degree rotation, but catching it with the cover
facing you. Why is this so hard?  You may wish to bring a book to the talk (either to flip or to read if the talk gets boring.)\\

\vspace{5mm} 
\huge   \textsc{Tuesday\\October 1, 2019 \\Hill 425 at 7:00 pm}
\\

\vspace{5mm}
\large
\emph{*Pizza and refreshments will be served}
\end{center}
\begin{center}
  \large  Join our email list! Email us at math.ruma@gmail.com, or
go to bit.ly/RUMA2019.\\
\end{center}



\end{document}

